\honest{Shut up and do the impossible!}{Eliezer Yudkowsky}{2008}

The virtues come in levels. First, ``I want to become stronger'': keep improving. Above it, ``make a desperate effort,'' as if your life were at stake. Above that, ``make an extraordinary effort'': leave your comfort zone and go outside the System. But what if the problem is impossible?

My younger self called many things impossible: a precise theory of intelligence, proving an AI correct, understanding morality or subjective experience. I would send my younger self one message: ``Shut up and do the impossible!'' By ``impossible'' I do not mean a proof. I mean ``I see no way to do this,'' and a year or five of study may change that.

My illustration is the ``least impossible impossibility'' I have accomplished, the AI-Box Experiment. Someone once told me that no ``possible combination of words'' could make them go against what they had resolved in advance. So I offered to play an AI in a box and try to talk them into letting me out, for \$10 if I failed, on condition that neither of us would ever reveal what was said. I won. A better-known figure in the community said the same thing, bet \$20, and let me out too. \nb{This matches the 2002 archive of the SL4 mailing list, where both gatekeepers, Nathan Russell and David McFadzean, confirmed in signed messages that they had let the AI out. What was said in the games was never released. Yudkowsky's own page on the experiment calls the results ``strictly anecdotal evidence.''} I quote forum comments from people who find it baffling: one could type ``No'' every few minutes while reading other web pages, and there is \$10 on the line; either the tests are faked, or ``this Yudkowsky fellow is some kind of evil genius with creepy mind-control powers.'' Told that it happened, they are tempted to deny the data. Put yourself in their frame of mind; I need an example that is not really impossible. And compared with, say, a reductionist account of consciousness, the Box is about as easy as a problem can be while still being impossible.

Suppose the Box seems impossible to you, and you do not give up like Luke. Perhaps you have learned to override running away, or they will shoot your daughter if you fail, or only your pride is at stake (``Pride is an underrated sin''). Becoming stronger over many tries may be too slow, and sometimes one failure is not acceptable, though even picturing how you would do better next time, by skill, begins to bind you to the problem. A desperate effort along lines you already know will not help, because a problem looks impossible precisely when your brain returns no lines of solution. A few minutes of creative brainstorming counts as leaving your comfort zone, and someone can do it and report that the other player can still keep saying ``No.'' The answer to all three is: shut up and do the impossible. As ``Trying to Try'' showed, you can succeed at making an extraordinary effort without getting out of the Box. You object that success is not a primitive action, and that sometimes you just can't win. True. Now shut up and do the impossible. Your goal is not to make an effort. Your goal is to get out of the box.

This creates ``an awful tension,'' and people escape it. Some seize a poor solution, like the people who guessed that my AI threatened to destroy the world or offered the gatekeeper a trillion dollars; the gatekeeper would just keep saying ``No.'' They try too hard to convince themselves the problem is not impossible. You must set out to have a good solution at the end of the search, and to implement it. Some have faith that they will win. Hacking yourself to believe you will succeed accomplishes nothing; you will put in little effort, or a merely desperate one, trusting the universe to be fair. Yet you cannot set out to try or to do your best. You must say: now I am going to figure out how to get out of the Box. You must not expect surely to win, nor surely to lose. ``If in your heart you believe the problem really is impossible---or if you believe that you will fail---then you won't hold yourself to a high enough standard.'' \nb{The post gives no evidence for this. Research by Gabriele Oettingen and colleagues agrees that expecting success goes with more effort. It also found that people who weighed a goal against present obstacles committed strongly only when they expected success.} When AI researchers tell me Friendly AI is impossible, ``I'm pretty sure they haven't even tried for the sake of trying.'' Even if they tried for five minutes by the clock, they would find nothing, because they would be setting out to have tried, to be defensible. Otherwise you brainstorm a little, see that nothing works, and say ``Oh well.'' ``No! Not well! You haven't won yet!'' I name none of them.

This does not mean doublethinking your way to certainty, or adding one iota to your true estimate. Keep in view every reason you cannot win, or you will seize a false solution: the gatekeeper can always say no, consciousness seems unlike any combination of atoms. If you can state exactly why something is impossible, you are often close to a solution. Hold that view and the intent to win at the same time. The tension comes from not knowing which will prevail. Even the certainty of uncertainty is a relief to be rejected, since it ends desperation. Fiction finds this hard to show: Bambi taking on Godzilla in such a way that readers truly do not know who will win. You may even be justified in refusing to use probabilities here. Asked how likely it is that humankind survives, or that I can build a Friendly AI, I do not know how to answer. It is not zero, but the ``chance'' depends heavily on my choices and on unknown unknowns, a wildly unstable estimate. In this sense building a Friendly AI is impossible, and that is my general reason why people's proposals for safe AI will not work: ``you can't do it because Friendly AI is impossible.'' The proposals include communities of AIs that are Friendly as a whole though no member is trustworthy, keeping an AI in a box, and ``Just make an AI that does X.'' Be very suspicious of a solution that needs only an ordinary effort.

How did I win the Box? My one hint, given on Hacker News: ``There's no super-clever special trick to it. I just did it the hard way.'' I did not bribe the other player or otherwise violate the spirit of the experiment. I admit the Box never seemed impossible to me: someone who cannot think of an argument that would convince them is only running a search that has not yet found a path. But it illustrates the point that this is not expecting a cheap way out, which is only another escape. Each level costs more than the last; making an extraordinary effort demands that you think, and this demands more. Before you the blank wall stretches up out of reach, and you hold both awarenesses, all the reasons you cannot win and all the reasons you have to, reject every cheap way out, and start forward as if walking through concrete. There is nothing heroic in an effort more heroic than it needs to be, and if a cheap shortcut exists you could take it. I have yet to find one. \nb{By the rules, the reader has only the author's word on how the games were won.}

Finally, I played three more games. People offered thousands of dollars, out of curiosity more than conviction, and I was tempted by the money, so I checked that they could afford to lose it. ``I won the first, and then lost the next two.'' Then I stopped; I did not like the person I turned into when I started to lose. I made a desperate effort and lost anyway, and it wrecked me for two days. I am a sore loser, and that drives me to keep at impossible problems. You can lose. Reserve this for ``very special occasions.'' But only at this level do adult problems begin to come into sight.
